\chapter{Justificativa e Histórico}
\label{chp:justificativa}

O \ppccurso da Universidade Federal de Uberlândia (UFU) é uma \textbf{reformulação e renomeação do curso de Engenharia de Controle e Automação}, já existente e ofertado pela Faculdade de Engenharia Elétrica (FEELT) (\autoref{sec:unidade_academica}), estruturada em conformidade com as Diretrizes Curriculares Nacionais do Curso de Graduação em Engenharia~\cite{cne_ces_2_2019,cne_ces_1_2021}. A reformulação parte da experiência consolidada da FEELT na formação em eletricidade, eletrônica, controle e automação — área em que a unidade já mantém corpo docente qualificado e infraestrutura laboratorial estabelecida por meio do curso ora reformulado — para ampliar essa formação em engenharia aplicada, incorporando robótica e inteligência artificial como ênfase curricular diferenciadora e passando a integrar, junto com o novo Curso Superior de Tecnologia em Automação Industrial, a Área Básica de Ingresso (ABI) descrita no \autoref{chp:abi}.

Esta reformulação foi viabilizada pela contemplação do curso pelo Programa Universidades Inovadoras e Sustentáveis~\cite{programa_universidades_inovadoras_sustentaveis}, iniciativa interministerial do Ministério da Educação (MEC) que redistribuiu novas vagas docentes e de técnicos de laboratório a universidades federais para a criação e a reformulação de cursos de graduação, com prioridade a áreas de tecnologia, inteligência artificial e STEM. Na FEELT, essa contemplação sustentou, conjuntamente: (i) a reformulação do curso de Engenharia de Controle e Automação, com mudança de oferta para o turno noturno e ampliação de \textbf{15 para 30} vagas ofertadas (\autoref{chp:identificacao}), passando a denominar-se Engenharia de Automação, Robótica e Inteligência Artificial; e (ii) a criação do novo Curso Superior de Tecnologia em Automação Industrial, também contemplado pelo mesmo programa para sua existência — justificando, em conjunto, a organização dos dois cursos em ABI (\autoref{chp:abi}; \autoref{sec:distincao_perfis}). \textcolor{red}{[CONFIRMAR FORMALMENTE JUNTO AO NDE/COLEGIADO, À FEELT E AOS CONSELHOS SUPERIORES DA UFU O ATO DE APROVAÇÃO DESTA REFORMULAÇÃO E RENOMEAÇÃO, BEM COMO O INSTRUMENTO ESPECÍFICO (EDITAL/PORTARIA) DE CONTEMPLAÇÃO PELO PROGRAMA UNIVERSIDADES INOVADORAS E SUSTENTÁVEIS]}

\section{Justificativa}

A expansão acelerada da automação industrial, impulsionada pela convergência entre eletrônica, controle, redes industriais, sistemas ciberfísicos, Internet Industrial das Coisas (IIoT) e inteligência artificial aplicada — o conjunto de transformações associado à Indústria 4.0 —, impôs uma demanda crescente por engenheiros(as) capazes de projetar, desenvolver, integrar e gerir sistemas automatizados industriais complexos. Ao mesmo tempo, o setor produtivo exige desses profissionais um perfil ético, comunicativo, colaborativo e sensível às realidades socioculturais em que estão inseridos.

O \ppccurso responde também ao chamado das políticas públicas e institucionais que estabelecem novos compromissos para os cursos de graduação, com destaque para a curricularização da extensão — historicamente associada à meta 12.7 do Plano Nacional de Educação 2014--2024~\cite{mec_pne_2014}, hoje regida pelo Plano Nacional de Educação 2026--2036~\cite{mec_pne_2026} — e, no âmbito da UFU, pela Resolução CONGRAD nº 177/2026~\cite{ufu_congrad_177_2026} (organização modular das Atividades Curriculares de Extensão detalhada em \autoref{sec:ace}). A Universidade Federal de Uberlândia, por meio do seu Plano Institucional de Desenvolvimento e Expansão (PIDE)~\cite{ufu_consun_31_2022}, orienta-se por essas diretrizes ao estimular a inovação pedagógica, a interdisciplinaridade, a equidade, a responsabilidade social e a formação orientada por competências, em consonância com as Diretrizes Curriculares Nacionais do Curso de Graduação em Engenharia~\cite{cne_ces_2_2019,cne_ces_1_2021}. A referência anterior à Resolução CNE/CP nº 1/2021, que trata das Diretrizes da Educação Profissional e Tecnológica, foi removida por não se aplicar a este Bacharelado em Engenharia.

A proposta também parte de um olhar atento para o contexto regional do Triângulo Mineiro e do interior do Brasil. A vocação tecnológica e industrial da região, os arranjos produtivos locais e a presença de indústrias com processos automatizados demandam uma formação versátil, aplicada e conectada às necessidades concretas do setor produtivo local, nacional e global — características centrais de um curso de engenharia orientado à indústria automatizada.

\subsection{Justificativa da Denominação e da Ênfase Curricular}

Diferentemente do Curso Superior de Tecnologia, cuja denominação decorre diretamente do Catálogo Nacional de Cursos Superiores de Tecnologia (CNCST), este curso não está vinculado a um catálogo de denominações fechado: as Diretrizes Curriculares Nacionais do Curso de Graduação em Engenharia~\cite{cne_ces_2_2019,cne_ces_1_2021} não fixam uma lista exaustiva de denominações, cabendo à instituição justificar a denominação proposta perante o MEC. A renomeação de Engenharia de Controle e Automação para Engenharia de Automação, Robótica e Inteligência Artificial (\autoref{sec:unidade_academica}) reflete a ampliação da ênfase curricular do curso reformulado: a denominação anterior — instituída pela Portaria MEC nº 1.694/1994, com atribuições profissionais consolidadas pela Resolução CONFEA nº 1.156/2025 (\autoref{sec:exercicio_profissional}) — permanece o precedente regulatório e institucional direto do curso, agora ampliado para refletir a incorporação de robótica e inteligência artificial como diferenciais curriculares. \textcolor{red}{[JUSTIFICAR FORMALMENTE PERANTE O NDE/COLEGIADO E O MEC A RENOMEAÇÃO PARA ``ENGENHARIA DE AUTOMAÇÃO, ROBÓTICA E INTELIGÊNCIA ARTIFICIAL'' NO PROCESSO DE REFORMULAÇÃO DO CURSO JÁ EXISTENTE — CONFIRMAR SE O MEC EXIGE NOVO ATO REGULATÓRIO ESPECÍFICO PARA A MUDANÇA DE DENOMINAÇÃO DE UM CURSO JÁ AUTORIZADO, OU SE BASTA A APROVAÇÃO PELOS CONSELHOS SUPERIORES DA UFU]}

A ênfase curricular em \ppcenfasecurricular\ (\autoref{chp:identificacao}) expressa a identidade pedagógica proposta neste PPC: além do núcleo central de eletricidade, eletrônica, instrumentação e controle, a matriz curricular deste curso incorpora um conjunto de conteúdos diferenciadores em robótica, sistemas inteligentes e inteligência artificial aplicada à automação. Essa ênfase é comunicada institucionalmente pelo nome de divulgação \ppccursomercadologico (ver \autoref{chp:identificacao}).

\subsection{Justificativa da Organização em Área Básica de Ingresso (ABI)}

O \ppccurso é organizado na ABI juntamente com o Curso Superior de Tecnologia em Automação Industrial (\autoref{chp:abi}), compartilhando os cinco primeiros períodos da trajetória acadêmica. Essa articulação se justifica pela proximidade das áreas tecnológicas envolvidas — automação industrial, robótica e inteligência artificial aplicada — e pela oportunidade de oferecer ao estudante uma base formativa comum antes de exigir a opção entre uma formação tecnológica de ciclo curto e diretamente aplicada (o Curso Superior de Tecnologia em Automação Industrial) e uma formação em Engenharia (o \ppccurso, \autoref{chp:justificativa}), de ciclo mais longo e com atribuições profissionais próprias.

A proposta de ABI busca ampliar a flexibilidade acadêmica sem descaracterizar as duas formações: o \ppccurso preserva identidade regulatória, matriz e perfil de egresso próprios (\autoref{chp:identificacao}; \autoref{chp:perfil_egresso}), assim como o Curso Superior de Tecnologia. A conclusão desse curso resulta em formação tecnológica completa e terminal. Qualquer continuidade posterior na Engenharia sem novo processo seletivo é apenas uma possibilidade a ser instituída e regulamentada pela UFU, conforme o \autoref{sec:permanencia}.

\subsection{Vagas Ofertadas}

O \ppccurso passa a ofertar \textbf{\ppcvagasofertadas}, no turno \ppcturno. A ampliação em relação às 15 vagas anteriormente ofertadas pelo curso de Engenharia de Controle e Automação, e a concomitante mudança de oferta do turno integral para o turno noturno, foram viabilizadas pelas novas vagas docentes e de técnicos de laboratório concedidas à FEELT no âmbito do Programa Universidades Inovadoras e Sustentáveis~\cite{programa_universidades_inovadoras_sustentaveis} (\autoref{sec:historico}) -- o mesmo programa que, simultaneamente, sustentou a criação do novo Curso Superior de Tecnologia em Automação Industrial, com o qual este curso passa a compartilhar a Área Básica de Ingresso (\autoref{chp:abi}).

\subsubsection{Continuidade, Ampliação e Base de Demanda}

Diferentemente do Curso Superior de Tecnologia, curso inteiramente novo, este curso é reformulação de oferta já consolidada: o antigo bacharelado em Engenharia de Controle e Automação já operava com 30 vagas autorizadas -- das quais 15 pelo SISU --, em regime integral e dez períodos~\cite{ufu_estudo_viabilidade_cst_2026}. A ampliação para 30 vagas via SISU e a migração para o turno noturno, portanto, não criam uma oferta a partir do zero: ampliam o acesso a uma formação cuja demanda e cuja base institucional -- corpo docente, laboratórios e produção científica do Departamento de Sistemas de Controle e Automação e do Laboratório de Automação e Sistemas Elétricos de Controle (LASEC) da FEELT~\cite{feelt_numeros} -- já estavam consolidadas.

A justificativa regional dessa ampliação apoia-se na mesma base econômica detalhada no PPC do Curso Superior de Tecnologia em Automação Industrial, com o qual este curso compartilha unidade acadêmica, corpo docente e infraestrutura: a Regional Vale do Paranaíba da FIEMG concentra 988.493 habitantes, valor adicionado industrial de R\$ 13,416 bilhões -- 78,8\% dele em Uberlândia --, 2.173 empresas de transformação (43.714 empregos, 31,3\% deles concentrados no setor de alimentos) e 82 empresas de SIUP (2.366 empregos)~\cite{fiemg_painel_vale_paranaiba_2026}, com âncoras industriais de processo contínuo e alta instrumentação como Cargill, Ambev, BRF e LD Celulose, esta última operando planta de 500 mil toneladas/ano de celulose solúvel no corredor Indianópolis--Araguari~\cite{cargill_uberlandia_investimento,ambev_uberlandia_investimento_2025,brf_uberlandia_unidade,ld_celulose_planta}. Esse parque industrial demanda não apenas a operação e a manutenção de sistemas automatizados -- campo de atuação do Curso Superior de Tecnologia --, mas também sua concepção, modelagem, projeto e integração em escala crescente de complexidade, incluindo sistemas ciberfísicos, robótica avançada, inteligência artificial aplicada e integração de plantas industriais completas -- campo próprio da atuação de um(a) \ppctitulacao.

\subsubsection{Capacidade Institucional}

A oferta apoia-se nas novas vagas docentes e de técnicos de laboratório concedidas à FEELT pelo Programa Universidades Inovadoras e Sustentáveis, somadas à estrutura já consolidada da unidade -- 1.373 estudantes, 72 docentes e 16 laboratórios em funcionamento~\cite{feelt_numeros}. \textcolor{red}{[INSERIR NÚMERO DE VAGAS DOCENTES EFETIVAMENTE CONCEDIDAS AO CURSO DE ENGENHARIA NO ÂMBITO DO PROGRAMA UNIVERSIDADES INOVADORAS E SUSTENTÁVEIS; INSERIR NÚMERO DE VAGAS DE TÉCNICOS DE LABORATÓRIO EFETIVAMENTE CONCEDIDAS; CONFIRMAR JUNTO AO NDE/COLEGIADO A CAPACIDADE NOMINAL DOS LABORATÓRIOS DESTINADOS ÀS ATIVIDADES PRÁTICAS DESTE CURSO, EM ESPECIAL OS COMPARTILHADOS COM O CURSO SUPERIOR DE TECNOLOGIA, E A EVENTUAL NECESSIDADE DE NOVAS AQUISIÇÕES]} A ampliação de 15 para 30 vagas no ingresso via SISU, ao lado da criação simultânea de 30 vagas noturnas no Curso Superior de Tecnologia, deve observar relação entre corpo docente/técnico e estudantes compatível com a natureza de projeto e de laboratório da formação em engenharia, especialmente em componentes de controle avançado, robótica e sistemas embarcados.

\subsubsection{Turno Noturno: Formação Científica sem Interrupção da Trajetória Profissional}

A migração para o turno noturno estende a um público mais amplo -- incluindo profissionais já inseridos na indústria regional que buscam formação em engenharia sem interromper sua atividade profissional -- o acesso a uma formação científica e tecnológica de maior profundidade: modelagem e análise de sistemas complexos, projeto e desenvolvimento de soluções de controle avançado, integração de tecnologias, pesquisa aplicada e inovação. Diferentemente do Curso Superior de Tecnologia, cuja terminalidade profissional é direta e cuja proposta central é a rápida inserção ou progressão do trabalhador-estudante, este curso preserva ciclo formativo mais longo e atribuições profissionais próprias da Engenharia (\autoref{sec:exercicio_profissional}) -- mas beneficia-se de argumento de fundo semelhante: um(a) estudante que permanece no ambiente industrial durante o dia pode trazer para os projetos, para o Trabalho de Conclusão de Curso e para as atividades de extensão problemas tecnológicos reais da indústria regional, em ciclo de aproximação contínua entre universidade e setor produtivo, sempre respeitados a confidencialidade, a propriedade intelectual das empresas e os regulamentos acadêmicos da UFU.

Da mesma forma que no Curso Superior de Tecnologia, é interesse deste curso compor parte de seu quadro docente efetivo -- via concurso público regular -- também com docentes em regime de 20 horas semanais, regime que, ao contrário da Dedicação Exclusiva (DE), não veda o exercício de outra atividade profissional remunerada (Lei nº 12.772/2012, art. 20)~\cite{lei_12772_2012,ufu_progep_regime_trabalho}. O objetivo é o mesmo: criar vínculo estável com profissionais que permanecem ativos na indústria regional durante o dia, trazendo experiência prática para a sala de aula e para as atividades de extensão, e aproximando os(as) estudantes das empresas da região -- mecanismo distinto da contratação, também possível, de professor substituto ou visitante de natureza temporária (Resolução CONSUN nº 24/1994)~\cite{ufu_consun_24_1994}. \textcolor{red}{[DEFINIR, JUNTO AO NDE E À DIREÇÃO DA FEELT, A PROPORÇÃO OU O NÚMERO DE VAGAS DOCENTES DESTE CURSO A SEREM DESTINADAS A REGIME DE 20 HORAS]}

\subsubsection{Complementaridade com o Curso Superior de Tecnologia}

A organização conjunta dos dois cursos em Área Básica de Ingresso (\autoref{chp:abi}) não decorre de equivalência entre eles, mas de complementaridade: o Curso Superior de Tecnologia forma para a integração, a operação, a manutenção e o \textit{troubleshooting} de sistemas automatizados já implantados; este curso forma para a concepção, o projeto, a modelagem e o desenvolvimento desses sistemas, incluindo seus aspectos mais avançados de controle, robótica e inteligência artificial. Compartilhar os cinco primeiros períodos da trajetória acadêmica, a infraestrutura laboratorial -- especialmente no período noturno -- e parte do corpo docente permite fortalecer os grupos de pesquisa e extensão da área de Controle e Automação da FEELT e ampliar, de forma conjunta, a interação com as empresas da região, sem que a racionalização de recursos se torne o argumento central dessa articulação: o que a justifica é a sinergia acadêmica entre dois perfis profissionais distintos e complementares, formados a partir da mesma base técnica.

\subsubsection{Transformação Digital da Indústria e Acompanhamento}

A ampliação de vagas neste curso também se justifica pela transformação tecnológica em curso na indústria regional e nacional: no âmbito da Nova Indústria Brasil, o MDIC estabeleceu a meta de digitalizar 50\% das empresas industriais brasileiras até 2033~\cite{mdic_nova_industria_brasil_2024}, e a Federação Internacional de Robótica registrou 542 mil robôs industriais instalados globalmente em 2024, mais do que o dobro de dez anos antes~\cite{abdi_boas_praticas_produtivas_2024} -- um movimento que demanda, além de operadores e integradores de sistemas automatizados, engenheiros(as) capazes de projetar e desenvolver as soluções de controle avançado, robótica, sistemas ciberfísicos e inteligência artificial industrial que sustentam essa transformação. O Núcleo Docente Estruturante e o Colegiado do Curso deverão acompanhar periodicamente indicadores de procura por vaga, ocupação, retenção, evasão, progressão acadêmica, inserção em estágio e empregabilidade dos egressos, em articulação com o acompanhamento equivalente do Curso Superior de Tecnologia, subsidiando eventual revisão do quantitativo de vagas com base em evidências.

\section{Histórico do Curso}
\label{sec:historico}

O \ppccurso encontra-se, no momento da elaboração deste PPC, em fase de proposta de reformulação (\texttt{perfil.status: proposta}, currículo \ppcversao), com Portaria de Autorização em tramitação no MEC (\autoref{chp:identificacao}). Trata-se da reformulação e renomeação do curso de Engenharia de Controle e Automação já existente e ofertado pela FEELT (\autoref{sec:unidade_academica}), e não de uma proposta de curso inteiramente novo — reformulação viabilizada pela contemplação do curso pelo Programa Universidades Inovadoras e Sustentáveis~\cite{programa_universidades_inovadoras_sustentaveis}.

\textcolor{red}{[CONFIRMAR O HISTÓRICO COMPLETO DO CURSO DE ENGENHARIA DE CONTROLE E AUTOMAÇÃO ORA REFORMULADO — ATO/RESOLUÇÃO DE CRIAÇÃO ORIGINAL; ATO/RESOLUÇÃO DE APROVAÇÃO DESTA REFORMULAÇÃO E RENOMEAÇÃO PELOS CONSELHOS SUPERIORES DA UFU; ANO PREVISTO DE IMPLANTAÇÃO DO CURRÍCULO REFORMULADO E DE INGRESSO DA PRIMEIRA TURMA NO NOVO TURNO NOTURNO]}

\section{Unidade Acadêmica}
\label{sec:unidade_academica}

O \ppccurso está vinculado à Faculdade de Engenharia Elétrica (FEELT) da Universidade Federal de Uberlândia (UFU), unidade acadêmica responsável por sua gestão pedagógica, administrativa e acadêmica. A FEELT atualmente abriga seis cursos de graduação presenciais — cinco localizados no campus Santa Mônica, em Uberlândia/MG (Engenharia Elétrica, Engenharia Biomédica, Engenharia de Computação, Engenharia Eletrônica e de Telecomunicações, Engenharia de Controle e Automação) e um no campus Patos de Minas/MG (Engenharia Eletrônica e de Telecomunicações) — além de dois programas de pós-graduação \textit{stricto sensu}: o Programa de Pós-Graduação em Engenharia Elétrica e o Programa de Pós-Graduação em Engenharia Biomédica.

Com a aprovação do Regimento Interno da FEELT~\cite{ufu_feelt_regimento_2023}, a estrutura organizacional da unidade está reorganizada em Departamentos Acadêmicos, cada um responsável por atividades de ensino, pesquisa e extensão em áreas específicas do conhecimento. Os departamentos atualmente constituídos na FEELT são:

\begin{itemize}
    \item Departamento de Sistemas de Potência e Energia (DSPE)
    \item Departamento de Engenharia Biomédica (DEB)
    \item Departamento de Computação e Inovação em Engenharia (DCIE)
    \item Departamento de Sistemas de Controle e Automação (DSCA)
    \item Departamento de Eletrônica e Telecomunicações (DETEL)
    \item Departamento de Máquinas Elétricas e Materiais (DMEM)
\end{itemize}

Pela natureza de sua matriz curricular — concentrada em eletricidade, eletrônica, instrumentação, controle, automação e redes industriais (\autoref{chp:identificacao}; aba \texttt{Nucleos} da matriz curricular) —, o \ppccurso tem forte interface com o Departamento de Sistemas de Controle e Automação (DSCA), com contribuição de docentes de outros departamentos da FEELT (DETEL, DMEM, DCIE) e de unidades acadêmicas externas, como a Faculdade de Matemática (FAMAT), o Instituto de Física (INFIS) e a Faculdade de Engenharia Química (FEQUI), na oferta de disciplinas obrigatórias (aba \texttt{Componentes} da matriz curricular). % REVISAR: unidade/departamento formalmente designado como responsável administrativo pelo curso ainda não confirmado nos documentos deste perfil.

Essa estrutura acadêmica consolidada garante ao curso uma base sólida de infraestrutura laboratorial e corpo docente qualificado para a formação de engenheiros(as) aptos(as) a atuar nos mais diversos contextos da automação, robótica e inteligência artificial aplicadas à indústria.

\section{Relação entre Sociedade e o Curso}

O \ppccurso orienta-se pela missão de formar profissionais capazes de responder às necessidades do setor produtivo, contribuindo para o avanço tecnológico, econômico e social da região e do país. Sua proposta articula-se com o contexto local e regional, buscando integrar-se ao ecossistema de inovação do Triângulo Mineiro.

A cidade de Uberlândia/MG, onde o curso está localizado, destaca-se como um dos principais polos tecnológicos e industriais do interior do Brasil, combinando infraestrutura logística estratégica, mão de obra qualificada e ambiente favorável ao desenvolvimento de negócios de base tecnológica. Diversas instituições e iniciativas locais vêm consolidando esse ecossistema, como a Secretaria Municipal de Desenvolvimento Econômico, Inovação e Turismo, a i9 Uberlândia, a Comunidade Colmeia, a Minas Startup, o Parque do Sabiá Tech e o SEBRAE/MG, fomentando a articulação entre universidades, indústrias, startups e o poder público.

A atuação de docentes da FEELT em projetos de pesquisa, desenvolvimento e inovação (PD\&I) tem resultado em parcerias com empresas e instituições dos setores elétrico e industrial — como ANEEL, CEMIG, Algar Telecom e ARCOM — e em projetos financiados por agências de fomento nacionais (CAPES, CNPq, FINEP, FAPEMIG), diretamente relacionadas ao campo de atuação de um(a) \ppctitulacao. % REVISAR: parcerias, convênios e indicadores de inserção profissional específicos deste curso (ainda em fase de proposta, sem egressos) não estão disponíveis; os dados acima refletem a atuação institucional da FEELT como um todo, não deste curso especificamente.

A proposta deste PPC reafirma o compromisso da UFU com a sociedade, respondendo à demanda por educação pública de qualidade voltada para uma área de alta relevância tecnológica e industrial, e promovendo uma formação alinhada às exigências contemporâneas da Indústria 4.0, da sustentabilidade e da inovação.
